\begin{sectionbox}
    \subsection{Frei ungedämpft}
    \textbf{Bewegungsgleichung}: $\ddot{x}+\omega_0^2x=0$\\
    $f_0=\frac{\omega_0}{2\pi} \qquad y_0=\frac{mg}{k} \qquad v_0 = \hat{x}\sqrt{\frac{k}{m}}$

    \subsubsection{Federpendel}
    $x(t)=x_{\ir max}\cos (\omega_0 t + \varphi_0) \qquad \omega_0 = \sqrt{\frac{k}{m}}$\\
    $x_{\ir max}=\sqrt{x_0^2+(\frac{v_0}{\omega_0})^2}$\\
    Bewegungsgleichung: $\ddot{x}+\frac{k}{m}x=0$

    \subsubsection{Fadenpendel}
    $x(t)=x_{\ir max}\cos (\omega_0 t + \varphi_0) \qquad \omega_0 = \sqrt\frac{g}{l}$\\
    $l=\frac{T^2g}{4\pi^2} \qquad h=l-l\cos \alpha$\\
    $x_{\ir max}=\sqrt{x_0^2+(\frac{v_0}{\omega_0})^2} \qquad E_{\ir ges} = mgl(\frac{3}{2}-\cos \alpha)$\\
    Bewegungsgleichung: $\ddot{\theta}+\frac{g}{l}\theta=0$

    \subsubsection{Torisionsschwingungen (Drehpendel) (TODO: s. Übung 9)}
    $\theta(t)=\theta_{\ir max}\cos(\omega_0 t + \varphi_0)$ \\
    $\omega_0 = \sqrt{\frac{k_t}{I}}=\sqrt{\frac{k_t}{mr^2}}$\\
    $I=\frac{T^2}{4\pi^2}k_t \qquad T_0=2\pi \sqrt{\frac{I}{k_t}}$\\
    Bewegungsgleichung: $\ddot{\theta}+\frac{k_t}{I}\theta=0$
\end{sectionbox}

\begin{sectionbox}
    \subsection{Frei gedämpft}
    \begin{tablebox}{ll}
        Stokesche Reibungskraft & $\vec F_R=-6 \pi \eta R \vec{v}=-2m\gamma\vec{v}$\\
        Resonanzfreq./gedämpfte Eigenfreq. & $\omega_D=\sqrt{\omega_0^2-\gamma^2}$ \\
        Bewegungsgleichung & $\ddot{x}+2\gamma \dot{x}+\omega_0^2x=0$
    \end{tablebox}
    $\gamma = \frac{\delta}{2m}=\frac{3\pi \eta R}{m} \quad \text{Reibungsterm/Abklingkoeffizient}, [\gamma] = \frac{1}{s}$ \\
    $\delta$: Dämpfungskonstante

    \subsubsection{Schwingfall (geringe Reibung) $\omega_0^2>\gamma^2$}
    $\omega_D=\sqrt{\omega_0^2-\gamma^2}<\omega_0$\\
    $x(t)=Ae^{-\gamma t}\cos(\omega_Dt)$

    \subsubsection{Aperiodischer Grenzfall $\omega_0^2=\gamma^2$}
    $\omega_D=0$\\
    $x(t)=Ae^{-\gamma t}  (1+\gamma t)$

    \subsubsection{Kriechfall (starke Reibung) $\omega_0^2<\gamma^2$}
    $\omega_D=i\gamma \pm i\sqrt{\gamma^2-\omega_0^2}$\\
    $x(t)=A_1e^{-\big(\gamma+\sqrt{\gamma^2-\omega_0^2}\big)t}+A_2e^{-\big(\gamma - \sqrt{\gamma^2-\omega_0^2}\big)t}$\\
\end{sectionbox}

\begin{sectionbox}
	\subsection{Erzwungen gedämpft}
    Gedämpfte Schwingung, die durch eine weitere Kraft dazu gezwungen wird, in Schwingung zu bleiben.
	Nach der Einschwungzeit soll das System mit der Kreisfrequenz $\Omega$ schwingen. \\

    \begin{tablebox}{ll}
        Resonanzfrequenz & $\Omega_{\ir res}=\sqrt{\omega_0^2-2\gamma^2}$ \\
        Bewegungsgleichung & $\ddot{x}+2\gamma \dot{x}+\omega_0^2x=\underbrace{F_0e^{-i\Omega t}}_{\text{externe Kraft}}$
    \end{tablebox}

    Auslenkung: $A = \frac{F_0}{m}\frac{1}{\sqrt{(\omega_0^2-\Omega^2)^2+4\gamma^2\Omega^2}}$ \\
    Phasenverschiebung: $\tan\varphi = \frac{2\gamma\Omega}{\omega_0^2-\Omega^2}$

    \begin{cookbox}{Resonanzfälle}
		\item $\Omega < \omega_0$ (Langsame Erregerfrequenz) \\
		$A\rightarrow \frac{F_0}{m\omega_0^2}=\frac{F_0}{k}$\qquad
		$\Delta \varphi = 0 \qquad \Omega \rightarrow 0$
		
        \item $\Omega \approx \omega_0$ (Katastrophe, Resonanz) für $\gamma \rightarrow 0$ \\
		$A_{\ir max}=\frac{F_0}{2m\gamma \omega_D} \rightarrow \frac{F_0}{\delta\omega_0} \rightarrow \infty$\\
		$\Omega_{\ir max}=\omega_D \rightarrow \omega_0 \qquad \Delta \varphi = \frac{\pi}{2}$
		
        \item $\Omega > \omega_0$ (Schnelle Erregerfrequenz) \\
		$\Delta \varphi = \pi \qquad \Omega \rightarrow \infty \qquad A\rightarrow 0$
	\end{cookbox}
    
\end{sectionbox}